\section{Intelligence per Watt: incorporar el consumo energético a la escalabilidad de capacidades}

\subsection{Del mainframe en la nube al local AI}

\videofigure{figures/intelligence-per-watt-title.jpg}{En la segunda mitad del curso se propone intelligence per watt, que mide la eficiencia del local AI mediante la capacidad del modelo disponible por unidad de consumo energético.}{00:40:08--00:42:00}

\videofigure{figures/mainframe-era.jpg}{La IA actual se asemeja a la era del mainframe: la mayor parte de la inferencia la proporciona infraestructura centralizada en la nube, mientras que la demanda crece más rápido que la oferta.}{00:43:48--00:45:22}

El curso observa que una gran cantidad de chatbot query reales solo requieren escritura, consulta de información, consejos generales y reasoning común, y no siempre necesitan un frontier cloud model. Si un modelo local de menos de aproximadamente 20B parámetros puede atender la mayoría de las solicitudes, es posible reducir la latencia de red, el costo en la nube y la exposición de la privacidad.

\videofigure{figures/local-model-progress.jpg}{La capacidad de los modelos locales pequeños crece rápidamente con cada generación de modelos y dispositivos, lo que crea condiciones para migrar parcialmente el cómputo de la nube al dispositivo.}{00:45:22--00:46:12}

\videofigure{figures/local-inference-question.jpg}{La pregunta central del sistema es: qué inferencias deben ejecutarse localmente y qué solicitudes vale la pena escalar a un frontier model en la nube.}{00:46:12--00:47:18}

\begin{importantbox}{Local/cloud no es una disyuntiva, sino un enrutamiento dinámico}
Las solicitudes simples, sensibles a la privacidad y de baja latencia se priorizan localmente; las que requieren el conocimiento de un modelo grande, un context largo o herramientas costosas se escalan a la nube. El propio enrutador debe estimar la dificultad, el riesgo y los recursos.
\end{importantbox}

\subsection{La definición de IPW}

\videofigure{figures/intelligence-per-watt-metric.jpg}{Intelligence per watt combina la proporción de solicitudes resolubles con el consumo energético promedio para comparar distintas combinaciones de modelo y hardware.}{00:47:18--00:48:18}

El indicador intuitivo usado en el curso puede escribirse como:

$$
\operatorname{IPW}(m,h)=
\frac{C(m)}{\overline P(m,h)},
$$

\begin{itemize}
    \item $m$: modelo;
    \item $h$: hardware de ejecución;
    \item $C(m)$: la proporción resoluble del modelo para el query workload objetivo o su puntuación de capacidad;
    \item $\overline P(m,h)$: el consumo energético promedio de dicho modelo en el hardware;
    \item cuanto mayor sea el IPW, más inteligencia útil se obtiene por unidad de consumo energético.
\end{itemize}

\videofigure{figures/efficiency-study.jpg}{El estudio mide la capacidad, la latencia, la precisión y el consumo energético en múltiples modelos locales, dispositivos y query workload reales.}{00:48:18--00:49:04}

\begin{warningbox}{El IPW debe vincularse al workload y a un umbral de calidad}
Si un modelo solo es competente en preguntas y respuestas simples, puede tener un consumo energético muy bajo pero ser incapaz de atender solicitudes complejas. Al comparar, es necesario fijar la distribución de consultas, el estándar de calidad de la salida, la longitud de context y la configuración de batch.
\end{warningbox}

\subsection{Una mejora de eficiencia de 5.3 veces en dos años}

\videofigure{figures/local-ai-findings.jpg}{Los modelos locales pueden atender con precisión alrededor del 88.7\% de las query del estudio; en dos años, la capacidad del modelo aumentó cerca de 3.1 veces, la eficiencia del hardware cerca de 1.7 veces, y el IPW en conjunto mejoró alrededor de 5.3 veces.}{00:49:04--00:52:18}

El curso también señala que el IPW del Apple M4 Max sigue siendo alrededor de 1.5 veces menor que el del B200, porque las GPU empresariales están optimizadas específicamente para el workload de LLM, mientras que los chips locales también deben atender gráficos, medios y aplicaciones de propósito general. Todavía existe un enorme margen de optimización futuro en los accelerator, kernel y arquitecturas de modelo del lado del dispositivo.

\begin{knowledgebox}{La aproximación de 5.3 veces proviene de la multiplicación de software y hardware}
$3.1\times1.7\approx5.27$. La compresión de modelos, los métodos de entrenamiento y la arquitectura elevan el numerator; el compilador, el kernel, el sistema de memoria y el chip elevan el denominator; la optimización conjunta suele ser más efectiva que optimizar solo un extremo.
\end{knowledgebox}

\subsection{Resumen de la sección}

Intelligence per watt reformula "un modelo más fuerte" como "ofrecer capacidad suficiente con menos energía en workload reales". El rápido progreso del local AI respalda la inferencia híbrida, pero la especialización del hardware y el enrutamiento de tareas siguen siendo claves.
